\scanpage{270}
\mspage{- 82 -}\hfill IX

par un tore. Nous avons posé comme d'habitude $A_n=A\times_S S_n$, où $S_n=\operatorname{Spec}(V/\underline m^{n+1})$. Nous nous intéresserons surtout au cas où $S$ est un trait, et où la fibre générique $A_K$ est un schéma abélien à réduction semi-stable sur $S$, $A$ étant le modèle de Néron de $A_K$.

Soit $T_{\unsure{o|0}}$ le tore maximal de $A_o$, de sorte que $B_o=A_o^o/T_o$ est un schéma abélien. La construction de 5.1 s'applique alors pour fournir un sous-tore formel $\hat T$ du complété formel $\hat A$. On sait que pour tout $n$, le quotient

\[
B_n=A_n^o/T_n\tag*{(7.1.1)}
\]

est représentable (SGA 3 $\mathrm{VI}_A$ 3.2), plus précisément c'est un schéma plat et de type fini sur $S_n$, dont l'image inverse sur $S_o$ est le schéma abélien $B_o$. Par suite $B_n$ est un \uline{schéma abélien} sur $S_n$, et pour $n$ variable, les $B_n$ se déduisent les uns des autres par changement de base, donc ils définissent un \uline{schéma abélien formel}

\[
\hat B=(B_n)_{n\geq0}\,.\tag*{(7.1.2)}
\]

Nous allons donner des conditions impliquant que ce schéma abélien est \uline{algébrisable} (EGA III 5.4.1), plus précisément que c'est le complété formel d'un schéma abélien projectif sur $S$, que nous noterons alors $B$. En vertu de loc.cit. $B$ est déterminé à isomorphisme unique près (comme $S$-schéma ; la structure de groupe de $\hat B$ provient alors de façon unique d'une structure de groupe de $B$, en vertu de loc. cit. également). Ce schéma abélien $B$ sera alors appelé \uline{la partie abélienne} du schéma $A$, et parfois noté $A^{\mathrm{ab}}$.

% archive sheet number: 269
% stamp: 400 (lower centre label; the pencilled 400 above it is also visible).
